\documentclass[11pt,a4paper]{article}
\usepackage[T1]{fontenc}
\usepackage[utf8]{inputenc}
\usepackage{lmodern,microtype,geometry}
\geometry{margin=27mm}
\usepackage{amsmath,amssymb,mathtools,amsthm,mathrsfs,booktabs,enumitem,xcolor,hyperref}
\hypersetup{colorlinks=true,linkcolor=blue!55!black,citecolor=blue!55!black,urlcolor=blue!55!black}
\setlist{nosep}
\newtheorem{theorem}{Theorem}[section]
\newtheorem{proposition}[theorem]{Proposition}
\newtheorem{corollary}[theorem]{Corollary}
\newtheorem{remark}[theorem]{Remark}
\newcommand{\Hil}{\mathcal H}
\newcommand{\BigHil}{\mathscr H}
\newcommand{\C}{\mathbb C}
\newcommand{\N}{\mathbb N}
\newcommand{\indop}{\operatorname{ind}}
\newcommand{\Span}{\operatorname{span}}
\newcommand{\diagop}{\operatorname{diag}}
\title{\textbf{Euler--Hilbert-Hotel Supersymmetry}\\\large Stable index classification, analytic zero modes, and binary-tetrahedral covariance}
\author{Adrian Lipa}
\date{22 September 2026}
\begin{document}
\maketitle
\begin{abstract}
We formulate two independent supersymmetry theorem modules. In the operator module, the standard Fredholm $\mathcal N=2$ factorization is combined with the Atiyah--J\"anich classification $\pi_0(\operatorname{Fred}(\mathcal H))\cong\mathbb Z$: the Witten index is the complete connected-component invariant, and unilateral shifts or adjoint shifts provide explicit Hilbert-Hotel normal representatives for every integer class. In the scalar Toeplitz sector,
\[
\Delta_W(T_f)=\operatorname{ind}(T_f)=-\operatorname{wind}(f),
\]
and the Toeplitz extension identifies this integer as the $K$-theory boundary image of the boundary-phase class. Finite Blaschke products sharpen the statement from an index count to an explicit zero-mode realization: $K_B=H^2\ominus BH^2$ is exactly the protected zero-mode space, and normalized reproducing-kernel overlaps recover the pseudohyperbolic geometry of the analytic zero locations.

Independently, the tetrahedral module identifies the regular qubit SIC as the projective Pauli $V_4$ orbit, extends the $Q_8$ spin lift to the binary tetrahedral group $2T$, and realizes a $3+1$ dimensional $\mathcal N=1$ supertranslation algebra after explicit CAR input. Equal tetrahedral weights define a timelike rest axis fixed by the tetrahedral rotational symmetry. Classical component theorems are explicitly marked as such; the contribution is the modular theorem packaging, exact normal forms, and explicit crosswalks between phase topology, Fredholm defects, protected zero modes, and finite spinorial covariance.
\end{abstract}
\section{Scope: two independent supersymmetry bricks}

The aim is deliberately modular. We isolate two exact supersymmetry constructions and keep their assumptions separate.

\begin{enumerate}[label=\textbf{S\arabic*.}]
\item An operator-index brick: Fredholm factorization, Toeplitz winding, Hilbert-Hotel shifts, and protected zero modes in $\mathcal N=2$ supersymmetric quantum mechanics.
\item A tetrahedral-spinor brick: a qubit SIC/Pauli orbit, its binary-tetrahedral spin symmetry, and an $\mathcal N=1$ supertranslation realization after explicit CAR input.
\end{enumerate}

No theorem requires the two bricks to be identified. Their value is that each can be checked, extended, or falsified independently.

\subsection{Euler/Pauli grading}

Euler's half-turn
\[
e^{i\pi}=-1
\]
appears in the spinorial lift as a central sign. For any Pauli axis,
\[
U_j(\pi)=\exp\!\left(-\frac{i\pi}{2}\sigma_j\right)=-i\sigma_j,
\qquad
U_j(\pi)^2=-I,
\qquad
U_j(\pi)^4=I.
\]
Thus an order-two projective half-turn lifts to an order-four spin operation.

The companion result SOH-G013 proves the determinant-one Pauli lifts
\[
\widetilde R=i\sigma_x,
\qquad
\widetilde E=i\sigma_z,
\qquad
\widetilde N=i\sigma_y,
\]
which generate the quaternion group $Q_8$ and project to the Klein four group $V_4$. The grading used below is
\[
\Gamma=\sigma_z\otimes I,
\qquad
\Gamma^2=I.
\]
The key strengthening in the present revision is that Euler phase also enters through winding: the same $U(1)$ phase variable that supplies the half-turn carries an integer topological charge on the boundary circle.

\section{Hilbert-Hotel shift and $\mathcal N=2$ supersymmetry}
Let $\Hil=\ell^2(\N)$ with orthonormal basis $e_1,e_2,\ldots$ and unilateral shift
\[
Se_n=e_{n+1}.
\]
Its adjoint obeys $S^\dagger e_1=0$ and $S^\dagger e_{n+1}=e_n$, hence
\begin{equation}\label{eq:shift}
S^\dagger S=I,\qquad SS^\dagger=I-P_1,\qquad P_1=|e_1\rangle\langle e_1|.
\end{equation}
This is the operator form of the Hilbert-Hotel self-embedding: the basis is isometrically mapped into the proper subset $\{e_2,e_3,\ldots\}$ with a rank-one defect.

\begin{proposition}
$S$ is Fredholm with $\indop(S)=-1$, and more generally $\indop(S^k)=-k$.
\end{proposition}
\begin{proof}
$S^\dagger S=I$ gives $\ker S=0$. The kernel of $S^\dagger$ is $\Span\{e_1\}$. For $S^k$, the cokernel is $\Span\{e_1,\ldots,e_k\}$.
\end{proof}

On $\BigHil=\Hil\oplus\Hil$ set
\[
\Gamma=\begin{pmatrix}I&0\\0&-I\end{pmatrix},\quad
\mathcal Q=\begin{pmatrix}0&0\\S&0\end{pmatrix},\quad
\mathcal Q^\dagger=\begin{pmatrix}0&S^\dagger\\0&0\end{pmatrix}.
\]

\begin{theorem}[Euler--Hilbert-Hotel supersymmetry]\label{thm:ehh}
The operators satisfy
\[
\mathcal Q^2=(\mathcal Q^\dagger)^2=0,
\qquad
\{\Gamma,\mathcal Q\}=\{\Gamma,\mathcal Q^\dagger\}=0,
\]
and
\[
H=\{\mathcal Q,\mathcal Q^\dagger\}
=\begin{pmatrix}I&0\\0&I-P_1\end{pmatrix}\ge0.
\]
For
\[
Q_1=\mathcal Q+\mathcal Q^\dagger,
\qquad
Q_2=i(\mathcal Q^\dagger-\mathcal Q),
\]
we have
\[
\boxed{\{Q_i,Q_j\}=2\delta_{ij}H},\qquad [H,Q_i]=0,\qquad \{\Gamma,Q_i\}=0.
\]
Thus the declared system is standard $\mathcal N=2$ supersymmetric quantum mechanics.
\end{theorem}
\begin{proof}
Nilpotence follows from the off-diagonal form. Direct multiplication gives $\mathcal Q^\dagger\mathcal Q=\diagop(S^\dagger S,0)$ and $\mathcal Q\mathcal Q^\dagger=\diagop(0,SS^\dagger)$. Equation~\eqref{eq:shift} then gives $H$. Expanding $Q_1,Q_2$ proves the remaining relations.
\end{proof}

\begin{corollary}[Witten/Fredholm index]
The positive sector has no zero mode and the negative sector has the protected zero mode $e_1$. With the displayed grading convention,
\[
\boxed{\Delta_W=-1=\indop(S)}.
\]
For the $k$-step shift,
\[
S^{\dagger k}S^k=I,\qquad S^kS^{\dagger k}=I-P_k,
\]
where $P_k$ projects onto $\Span\{e_1,\ldots,e_k\}$. Therefore
\[
\boxed{\Delta_W(H_k)=-k=\indop(S^k)},
\]
and exactly $k$ zero modes are protected in the negative grading sector.
\end{corollary}

\subsection{Zero-centred analytic carrier}
On Hardy space $H^2(\mathbb D)$, multiplication by $z$ is exactly the unilateral shift on the basis $1,z,z^2,\ldots$. The distinguished zero $z=0$ determines the codimension-one subspace
\[
zH^2(\mathbb D)=\{f\in H^2(\mathbb D):f(0)=0\}.
\]
Hence a zero-centred complex analytic domain supports the theorem once the Hilbert carrier is specified. The point $0$ alone does not force a unique supercharge.

\section{Tetrahedral null-generated causal cone}
Define
\[
\mathbf n_1=\frac{(1,1,1)}{\sqrt3},\quad
\mathbf n_2=\frac{(1,-1,-1)}{\sqrt3},\quad
\mathbf n_3=\frac{(-1,1,-1)}{\sqrt3},\quad
\mathbf n_4=\frac{(-1,-1,1)}{\sqrt3}.
\]
Then
\[
\sum_a\mathbf n_a=0,
\qquad
\mathbf n_a\cdot\mathbf n_b=-\frac13\ (a\ne b),
\qquad
\frac14\sum_a\mathbf n_a\mathbf n_a^T=\frac13I_3.
\]
With Minkowski metric $\eta=\diagop(+,-,-,-)$ define $k_a=(1,\mathbf n_a)$. Each $k_a$ is future null, and
\[
\frac14\sum_a k_a=(1,0,0,0).
\]

For $p_a\ge0$, set $P=\sum_ap_ak_a$.
\begin{theorem}[Exact tetrahedral causal norm]\label{thm:cone}
\[
\boxed{P^2=\frac83\sum_{a<b}p_ap_b\ge0.}
\]
Thus the positive span of the four null rays is contained in the future Lorentz cone, and $P$ is null exactly when at most one $p_a$ is nonzero.
\end{theorem}
\begin{proof}
$k_a^2=0$, while for $a\ne b$,
\[
k_a\cdot k_b=1-\mathbf n_a\cdot\mathbf n_b=\frac43.
\]
Expanding $P^2$ gives the formula.
\end{proof}
At fixed $P^0=1$ the section is the regular tetrahedron inside the unit ball. Its four vertices are null directions; nontrivial convex mixtures are timelike. It is therefore a null-generated tetrahedral causal cone, not a replacement for the full null boundary.

\section{Spinorial and supersymmetric lift}
Let $\sigma^\mu=(I_2,\sigma_x,\sigma_y,\sigma_z)$ and
\[
K_a=I_2+\mathbf n_a\cdot\boldsymbol\sigma.
\]
Since $|\mathbf n_a|=1$, $K_a$ is positive semidefinite, rank one, and $\det K_a=0$. Hence there is a two-spinor $\lambda_a\in\C^2$ such that
\[
\boxed{K_a=\lambda_a\lambda_a^\dagger}.
\]
The factorization has the full phase gauge $\lambda_a\sim e^{i\chi_a}\lambda_a$; the sign $\lambda_a\mapsto-\lambda_a$ is the Euler half-turn element, not the entire gauge freedom.

The geometry supplies the null spinors but not fermionic operators. Attach four modes $f_a,f_a^\dagger$ satisfying
\[
\{f_a,f_b^\dagger\}=\delta_{ab},
\qquad
\{f_a,f_b\}=\{f_a^\dagger,f_b^\dagger\}=0.
\]
For $p_a\ge0$ define
\[
P_{\alpha\dot\beta}=\sum_a p_a\lambda_{a\alpha}\bar\lambda_{a\dot\beta},
\]
\[
Q_\alpha=\sqrt2\sum_a\sqrt{p_a}\lambda_{a\alpha}f_a^\dagger,
\qquad
\bar Q_{\dot\beta}=\sqrt2\sum_a\sqrt{p_a}\bar\lambda_{a\dot\beta}f_a.
\]

\begin{theorem}[Tetrahedral $\mathcal N=1$ supertranslation algebra]\label{thm:tetra-susy}
\[
\boxed{\{Q_\alpha,\bar Q_{\dot\beta}\}=2P_{\alpha\dot\beta}},
\qquad
\{Q_\alpha,Q_\beta\}=0,
\qquad
\{\bar Q_{\dot\alpha},\bar Q_{\dot\beta}\}=0.
\]
Thus the tetrahedral momentum cone admits an exact $3+1$ dimensional $\mathcal N=1$ supertranslation realization once one CAR mode is supplied for each null ray.
\end{theorem}
\begin{proof}
Using the CAR,
\[
\{Q_\alpha,\bar Q_{\dot\beta}\}
=2\sum_{a,b}\sqrt{p_ap_b}\lambda_{a\alpha}\bar\lambda_{b\dot\beta}\{f_a^\dagger,f_b\}
=2\sum_ap_a\lambda_{a\alpha}\bar\lambda_{a\dot\beta}.
\]
The same-chirality anticommutators vanish because the creation operators mutually anticommute, as do the annihilation operators.
\end{proof}

Theorem~\ref{thm:tetra-susy} proves that the tetrahedral spinorial carrier admits a natural supersymmetric lift after CAR input. It does not prove that geometry alone derives physical fermions.

\section{Relation of the two constructions}
Theorem~\ref{thm:ehh} is a $0+1$ dimensional supersymmetric quantum-mechanical construction whose invariant is a Fredholm/Witten index. Theorem~\ref{thm:tetra-susy} is a $3+1$ dimensional supertranslation construction on a discrete null-generated momentum cone. They share a binary grading and a factorization of a positive object as a square or anticommutator, but no functorial identification between the two sectors is proved here.

\section{Nested tetrahedral sections and Euler's Infinity}
One may consider nested tetrahedral sections
\[
\mathcal T_n=r_n\operatorname{conv}\{\mathbf n_1,\ldots,\mathbf n_4\},
\qquad r_n\to0,
\]
around the fixed timelike axis $(1,0,0,0)$. This provides a mathematically clean carrier for an ``Euler Infinity'' programme. However, the scale law $r_n$ must be derived from a declared dynamics. No specific law, including $r_n=e^{-n}$, follows from Euler's identity alone.

The exact supersymmetry theorems above do not require an infinite nest. A canonical inter-section evolution and its compatibility with the spinor/supercharge structures remain \textbf{OPEN / MODEL-LEVEL}.

\section{Proof firewall and conclusion}
The exact results are:
\begin{enumerate}[label=\textbf{E\arabic*.}]
\item the unilateral-shift $\mathcal N=2$ algebra;
\item $\Delta_W(H_k)=\indop(S^k)=-k$ with the declared grading convention;
\item the regular-tetrahedron isotropy and causal-cone identity of Theorem~\ref{thm:cone};
\item rank-one spinor factorization of the four null generators;
\item the $\mathcal N=1$ supertranslation algebra after explicit CAR modes are supplied.
\end{enumerate}

Not established are: physical supersymmetry of Nature; a theory of everything; derivation of the CAR modes from tetrahedral geometry alone; a unique supercharge from a bare pointed complex plane; a canonical scale law for the infinite tetrahedral nest; or any consequence for RH, Collatz, Twin Prime, or another open problem.

The compact operator result is
\[
\boxed{
\text{unilateral shift}+\mathbb Z_2\text{ grading}
\Longrightarrow \mathcal N=2\text{ SUSY},
\qquad
\Delta_W=\indop(S^k)=-k.}
\]
The Hilbert-Hotel defect is therefore literally a supersymmetric index in the declared operator model.

Independently, the tetrahedral spinor/CAR construction gives
\[
\boxed{
\{Q_\alpha,\bar Q_{\dot\beta}\}
=2\sigma^\mu_{\alpha\dot\beta}P_\mu}
\]
on the tetrahedral momentum cone. This upgrades the tetrahedral spinorial picture from analogy to an exact algebraic supersymmetry realization while preserving the crucial boundary that the fermionic degrees of freedom are supplied rather than derived.

\begin{thebibliography}{99}
\bibitem{Witten1981} E. Witten, ``Dynamical Breaking of Supersymmetry,'' \emph{Nuclear Physics B} \textbf{188} (1981), 513--554.
\bibitem{Cooper1995} F. Cooper, A. Khare, and U. Sukhatme, ``Supersymmetry and quantum mechanics,'' \emph{Physics Reports} \textbf{251} (1995), 267--385.
\bibitem{Janich1965} K. J\"anich, ``Vektorraumb\"undel und der Raum der Fredholm-Operatoren,'' \emph{Mathematische Annalen} \textbf{161} (1965), 129--142, doi:10.1007/BF01360851.

\bibitem{Douglas1998} R. G. Douglas, \emph{Banach Algebra Techniques in Operator Theory}, 2nd ed., Graduate Texts in Mathematics 179, Springer, 1998.

\bibitem{Sarason1994} D. Sarason, \emph{Sub-Hardy Hilbert Spaces in the Unit Disk}, Wiley, 1994.

\bibitem{BaumVanErp2021} P. F. Baum and E. van Erp, ``The Index Theorem for Toeplitz Operators as a Corollary of Bott Periodicity,'' \emph{Quarterly Journal of Mathematics} \textbf{72} (2021), 547--569, doi:10.1093/qmath/haab008.

\bibitem{Izumi2026} M. Izumi, ``A generalized winding number formula for the Witten index of a Toeplitz operator,'' \emph{Journal of Spectral Theory}, online first (2026), doi:10.4171/JST/593.

\bibitem{Renes2004} J. M. Renes, R. Blume-Kohout, A. J. Scott, and C. M. Caves, ``Symmetric informationally complete quantum measurements,'' \emph{Journal of Mathematical Physics} \textbf{45} (2004), 2171--2180, doi:10.1063/1.1737053.

\bibitem{ConwaySmith} J. H. Conway and D. A. Smith, \emph{On Quaternions and Octonions}, A K Peters, 2003.

\bibitem{PenroseRindler} R. Penrose and W. Rindler, \emph{Spinors and Space-Time, Vol. 1}, Cambridge University Press, 1984.
\bibitem{LipaInfinities} A. Lipa, \emph{Infinities}, research repository, 2026. \url{https://github.com/AdrianLipa90/Infinities}
\bibitem{LipaSOH} A. Lipa, ``SOH-G013 -- Pauli / SU(2) Spinor Lift of the Projective Operator Algebra,'' \emph{Secret-of-a-Half}, 2026. \url{https://github.com/AdrianLipa90/secret-of-a-half/blob/main/research/SOH_G013_PAULI_SPINOR_LIFT_V0_1.md}
\end{thebibliography}
\end{document}
